\documentclass[12pt,a4paper]{article}

\usepackage[T2A]{fontenc}
\usepackage[utf8]{inputenc}
\usepackage[russian]{babel}
\usepackage{amsmath,amssymb,amsthm,mathtools}
\usepackage{enumitem}
\usepackage{geometry}
\usepackage{array}
\usepackage{booktabs}
\usepackage{tikz}
\usepackage{listings}
\usepackage{xcolor}

\geometry{left=25mm,right=20mm,top=20mm,bottom=20mm}

\theoremstyle{definition}
\newtheorem{definition}{Определение}
\newtheorem{example}{Пример}
\newtheorem{remark}{Замечание}

\theoremstyle{plain}
\newtheorem{theorem}{Теорема}
\newtheorem{proposition}{Утверждение}

\lstdefinestyle{pseudo}{
    basicstyle=\ttfamily\small,
    keywordstyle=\bfseries\color{blue!60!black},
    commentstyle=\itshape\color{gray!70!black},
    frame=single,
    breaklines=true,
    columns=fullflexible,
    keepspaces=true,
    showstringspaces=false
}

\title{Билет №50 \\ \small Понятие о вычислительном эксперименте и его инструментальной поддержке. (СпБПУ)}
\author{Сериков Владислав Александрович}
\date{16 мая 2026}

\begin{document}

\maketitle
\tableofcontents
\newpage

\section{Введение}

Вычислительный эксперимент является одной из центральных форм современной научной и инженерной деятельности. Он занимает промежуточное положение между теоретическим анализом и натурным экспериментом: исследуемый объект заменяется математической моделью, а свойства этой модели изучаются с помощью численных методов и программной реализации на вычислительной технике.

Вычислительный эксперимент особенно важен в тех случаях, когда:
\begin{itemize}
    \item натурный эксперимент слишком дорог, опасен или невозможен;
    \item объект имеет большие пространственные или временные масштабы;
    \item требуется многократно исследовать поведение системы при различных параметрах;
    \item необходимо прогнозирование;
    \item нужно проверить гипотезы до проведения реального эксперимента;
    \item требуется оптимизация конструкции, технологии или режима управления.
\end{itemize}

Примеры областей применения:
\begin{itemize}
    \item гидро- и газодинамика;
    \item теплопроводность и диффузия;
    \item механика деформируемого твёрдого тела;
    \item климатическое моделирование;
    \item плазменная физика;
    \item биология и медицина;
    \item экономика;
    \item социальные и транспортные системы;
    \item машинное обучение и анализ данных.
\end{itemize}

\section{Математическое моделирование как основа вычислительного эксперимента}

\begin{definition}
\textbf{Модель} --- это объект любой природы, который в определённом отношении замещает исследуемый объект и позволяет получить информацию о нём.
\end{definition}

\begin{definition}
\textbf{Математическая модель} --- это описание объекта, процесса или явления с помощью математических понятий: переменных, функций, уравнений, неравенств, операторов, вероятностных распределений, алгоритмических правил и т.д.
\end{definition}

Математическая модель обычно включает:
\begin{itemize}
    \item входные данные;
    \item искомые величины;
    \item параметры модели;
    \item уравнения или алгоритмические правила;
    \item начальные условия;
    \item граничные условия;
    \item критерии качества или целевую функцию;
    \item допущения и область применимости.
\end{itemize}

\begin{definition}
\textbf{Объект моделирования} --- реальная или абстрактная система, свойства которой изучаются.

\textbf{Модельный объект} --- математическая конструкция, используемая вместо исходного объекта.
\end{definition}

\subsection{Этапы математического моделирования}

Типичная схема математического моделирования имеет вид:
\[
\text{объект} \longrightarrow
\text{содержательная модель} \longrightarrow
\text{математическая модель} \longrightarrow
\text{численная модель} \longrightarrow
\text{программа} \longrightarrow
\text{результаты}.
\]

\begin{center}
\begin{tikzpicture}[node distance=2.2cm, >=latex]
\node[draw, rounded corners, align=center] (obj) {Реальный\\объект};
\node[draw, rounded corners, right of=obj, align=center] (phys) {Содержательная\\модель};
\node[draw, rounded corners, right of=phys, align=center] (math) {Математическая\\модель};
\node[draw, rounded corners, below of=math, align=center] (num) {Численная\\модель};
\node[draw, rounded corners, left of=num, align=center] (prog) {Программная\\реализация};
\node[draw, rounded corners, left of=prog, align=center] (res) {Анализ\\результатов};

\draw[->] (obj) -- (phys);
\draw[->] (phys) -- (math);
\draw[->] (math) -- (num);
\draw[->] (num) -- (prog);
\draw[->] (prog) -- (res);
\draw[->, dashed] (res) to[bend left=35] (obj);
\draw[->, dashed] (res) to[bend right=35] (math);
\end{tikzpicture}
\end{center}

Пунктирные стрелки подчёркивают итерационный характер моделирования: результаты вычислений могут приводить к уточнению модели, параметров, численного метода или исходных гипотез.

\section{Понятие вычислительного эксперимента}

\begin{definition}
\textbf{Вычислительный эксперимент} --- это целенаправленное исследование математической модели изучаемого объекта с помощью вычислительных алгоритмов и компьютерной реализации с последующим анализом, интерпретацией и проверкой полученных результатов.
\end{definition}

Иными словами, вычислительный эксперимент представляет собой эксперимент не непосредственно над физическим объектом, а над его математической и программной моделью.

\subsection{Структура вычислительного эксперимента}

Вычислительный эксперимент включает следующие основные компоненты:
\begin{enumerate}
    \item постановку прикладной задачи;
    \item построение математической модели;
    \item анализ корректности модели;
    \item выбор численного метода;
    \item построение дискретной модели;
    \item программную реализацию;
    \item верификацию программы;
    \item проведение серии расчётов;
    \item анализ чувствительности и устойчивости;
    \item сравнение с теорией, экспериментом или эталонными решениями;
    \item интерпретацию результатов.
\end{enumerate}

\subsection{Отличие вычислительного эксперимента от натурного}

\begin{center}
\begin{tabular}{p{0.42\textwidth} p{0.42\textwidth}}
\toprule
\textbf{Натурный эксперимент} & \textbf{Вычислительный эксперимент} \\
\midrule
Проводится над реальным объектом & Проводится над математической и программной моделью \\
Требует лабораторной или промышленной установки & Требует вычислительной среды и программного обеспечения \\
Может быть дорогим или опасным & Обычно дешевле и безопаснее \\
Ограничен измерительными возможностями & Позволяет наблюдать недоступные в реальности величины \\
Содержит физические погрешности измерений & Содержит модельные, численные и программные ошибки \\
Не всегда воспроизводим & Воспроизводим при фиксации кода, данных и среды \\
\bottomrule
\end{tabular}
\end{center}

\subsection{Отличие вычислительного эксперимента от аналитического исследования}

Аналитическое исследование стремится получить решение в явной форме, например:
\[
u(x,t) = f(x,t).
\]

В вычислительном эксперименте решение обычно получается приближённо:
\[
u(x_i,t_n) \approx u_i^n,
\]
где $x_i$ и $t_n$ --- узлы сетки по пространству и времени.

Аналитические методы дают точные соотношения, но применимы к сравнительно узкому классу задач. Вычислительные методы позволяют исследовать сложные нелинейные, многомерные, нестационарные системы, для которых точное решение недоступно.

\section{Основные этапы вычислительного эксперимента}

\subsection{Постановка задачи}

На первом этапе формулируется содержательная постановка:
\begin{itemize}
    \item что требуется найти;
    \item какие величины известны;
    \item какие параметры изменяются;
    \item какие ограничения существенны;
    \item какие результаты считаются приемлемыми.
\end{itemize}

Например, для задачи теплопроводности требуется определить распределение температуры $u(x,t)$ в стержне при заданных начальном и граничных условиях.

\subsection{Построение математической модели}

Математическая модель может иметь вид:
\begin{itemize}
    \item алгебраических уравнений;
    \item обыкновенных дифференциальных уравнений;
    \item уравнений в частных производных;
    \item интегральных уравнений;
    \item стохастических моделей;
    \item дискретных динамических систем;
    \item агентных моделей;
    \item оптимизационных задач.
\end{itemize}

\begin{example}
Модель теплопроводности в одномерном стержне:
\[
\frac{\partial u}{\partial t}
=
a^2 \frac{\partial^2 u}{\partial x^2},
\qquad
0 < x < L,\quad t>0,
\]
с начальными и граничными условиями:
\[
u(x,0)=\varphi(x),
\]
\[
u(0,t)=\mu_0(t), \qquad u(L,t)=\mu_L(t).
\]
Здесь $u(x,t)$ --- температура, $a^2$ --- коэффициент температуропроводности.
\end{example}

\subsection{Анализ корректности}

\begin{definition}
Задача называется \textbf{корректно поставленной по Адамару}, если выполнены три условия:
\begin{enumerate}
    \item решение существует;
    \item решение единственно;
    \item решение непрерывно зависит от исходных данных.
\end{enumerate}
\end{definition}

Непрерывная зависимость от данных означает, что малые изменения входных данных вызывают малые изменения решения.

Если задача некорректна, вычислительный эксперимент может давать неустойчивые и физически бессмысленные результаты.

\subsection{Дискретизация}

Для численного решения непрерывная задача заменяется дискретной.

Например, вводится сетка:
\[
x_i = ih,\qquad i=0,1,\dots,N,\qquad h=\frac{L}{N},
\]
\[
t_n = n\tau,\qquad n=0,1,\dots,M.
\]

Значение функции $u(x_i,t_n)$ приближается сеточной функцией:
\[
u(x_i,t_n) \approx u_i^n.
\]

Производные заменяются конечными разностями:
\[
\frac{\partial u}{\partial t}(x_i,t_n)
\approx
\frac{u_i^{n+1}-u_i^n}{\tau},
\]
\[
\frac{\partial^2 u}{\partial x^2}(x_i,t_n)
\approx
\frac{u_{i+1}^n - 2u_i^n + u_{i-1}^n}{h^2}.
\]

Тогда уравнение теплопроводности аппроксимируется явной схемой:
\[
\frac{u_i^{n+1}-u_i^n}{\tau}
=
a^2
\frac{u_{i+1}^n - 2u_i^n + u_{i-1}^n}{h^2}.
\]

Отсюда:
\[
u_i^{n+1}
=
u_i^n + r\left(u_{i+1}^n - 2u_i^n + u_{i-1}^n\right),
\qquad
r=\frac{a^2\tau}{h^2}.
\]

\subsection{Построение алгоритма}

\begin{definition}
\textbf{Алгоритм} --- точное описание конечной последовательности действий, приводящей от исходных данных к результату.
\end{definition}

Для вычислительного эксперимента алгоритм должен быть:
\begin{itemize}
    \item однозначным;
    \item конечным;
    \item реализуемым на компьютере;
    \item устойчивым к малым ошибкам данных и округления;
    \item достаточно эффективным по времени и памяти.
\end{itemize}

\subsection{Программная реализация}

На данном этапе алгоритм реализуется в виде программы. Важно обеспечить:
\begin{itemize}
    \item правильную структуру кода;
    \item контроль входных данных;
    \item обработку исключительных ситуаций;
    \item тестирование;
    \item документирование;
    \item воспроизводимость вычислений.
\end{itemize}

\subsection{Верификация и валидация}

\begin{definition}
\textbf{Верификация} --- проверка того, что программа правильно решает выбранную математическую и численную задачу.
\end{definition}

\begin{definition}
\textbf{Валидация} --- проверка того, что математическая модель адекватно описывает реальный объект в рассматриваемой области применимости.
\end{definition}

Кратко:
\[
\text{верификация: ``правильно ли решаем уравнения?''}
\]
\[
\text{валидация: ``правильные ли уравнения решаем?''}
\]

\section{Ошибки вычислительного эксперимента}

Результат вычислительного эксперимента всегда содержит ошибки различной природы.

\subsection{Основные типы ошибок}

\begin{enumerate}
    \item \textbf{Модельная ошибка} --- связана с упрощениями и допущениями при построении математической модели.
    \item \textbf{Ошибка исходных данных} --- возникает из-за неточности измерений и параметров.
    \item \textbf{Ошибка аппроксимации} --- возникает при замене непрерывной задачи дискретной.
    \item \textbf{Итерационная ошибка} --- появляется при приближённом решении алгебраических систем или нелинейных уравнений.
    \item \textbf{Ошибка округления} --- связана с конечной разрядностью компьютерной арифметики.
    \item \textbf{Программная ошибка} --- ошибка реализации алгоритма.
    \item \textbf{Интерпретационная ошибка} --- неправильное понимание численного результата.
\end{enumerate}

Общую ошибку можно условно представить как:
\[
E_{\text{total}}
=
E_{\text{model}}
+
E_{\text{data}}
+
E_{\text{approx}}
+
E_{\text{iter}}
+
E_{\text{round}}
+
E_{\text{prog}}.
\]

Такое разложение является схематическим, поскольку в реальности ошибки могут взаимодействовать нелинейно.

\subsection{Абсолютная и относительная погрешности}

\begin{definition}
Пусть $x$ --- точное значение, $\widetilde{x}$ --- приближённое значение. Тогда
\[
\Delta x = |\widetilde{x}-x|
\]
называется \textbf{абсолютной погрешностью}.
\end{definition}

\begin{definition}
Если $x\neq 0$, то величина
\[
\delta x = \frac{|\widetilde{x}-x|}{|x|}
\]
называется \textbf{относительной погрешностью}.
\end{definition}

\begin{example}
Пусть точное значение равно $x=100$, а приближённое $\widetilde{x}=98$.
Тогда
\[
\Delta x = |98-100|=2,
\]
\[
\delta x = \frac{2}{100}=0.02=2\%.
\]
\end{example}

\section{Аппроксимация, устойчивость и сходимость}

\subsection{Аппроксимация}

Рассмотрим операторное уравнение:
\[
Lu=f.
\]

Пусть $L_h$ --- дискретный оператор, а $u_h$ --- сеточная функция. Разностная схема имеет вид:
\[
L_h u_h = f_h.
\]

\begin{definition}
Разностная схема \textbf{аппроксимирует} исходную дифференциальную задачу, если при подстановке в схему точного решения $u$ невязка стремится к нулю при $h\to 0$:
\[
\|L_h [u] - f_h\| \to 0,
\qquad h\to 0.
\]
\end{definition}

Если
\[
\|L_h [u] - f_h\| \le C h^p,
\]
то говорят, что схема имеет порядок аппроксимации $p$ по параметру $h$.

\subsection{Устойчивость}

\begin{definition}
Разностная схема называется \textbf{устойчивой}, если малые возмущения входных данных, правой части, начальных и граничных условий вызывают равномерно малые изменения численного решения.
\end{definition}

В простейшей операторной форме устойчивость означает существование постоянной $C$, не зависящей от шага сетки $h$, такой что:
\[
\|u_h\| \le C \|f_h\|.
\]

\subsection{Сходимость}

\begin{definition}
Разностная схема называется \textbf{сходящейся}, если её решение $u_h$ стремится к точному решению $u$ исходной задачи при измельчении сетки:
\[
\|u_h - [u]\| \to 0,
\qquad h\to 0.
\]
\end{definition}

\subsection{Теорема Лакса--Рихтмайера}

\begin{theorem}[Теорема Лакса--Рихтмайера об эквивалентности]
Для корректно поставленной линейной начальной задачи и её согласованной разностной аппроксимации устойчивость разностной схемы эквивалентна её сходимости.
\end{theorem}

Иначе говоря, для линейных корректных задач:
\[
\text{аппроксимация} + \text{устойчивость}
\quad \Longleftrightarrow \quad
\text{сходимость}.
\]

\begin{proof}
Приведём доказательство в стандартной операторной форме.

Пусть точная задача имеет вид
\[
Lu=f,
\]
а разностная схема:
\[
L_h u_h = f_h.
\]

Обозначим через $[u]$ проекцию точного решения на сетку. Ошибка численного решения:
\[
e_h = u_h - [u].
\]

Подставим $u_h$ и $[u]$ в разностный оператор:
\[
L_h u_h = f_h.
\]

Для точного решения, подставленного в схему, вообще говоря,
\[
L_h [u] = f_h + \psi_h,
\]
где $\psi_h$ --- ошибка аппроксимации, или невязка.

Тогда
\[
L_h u_h - L_h [u] = f_h - (f_h+\psi_h),
\]
то есть
\[
L_h (u_h-[u]) = -\psi_h.
\]

Следовательно,
\[
L_h e_h = -\psi_h.
\]

Если схема устойчива, то существует постоянная $C$, не зависящая от $h$, такая что
\[
\|e_h\| \le C\|\psi_h\|.
\]

Так как схема согласована, то есть аппроксимирует исходную задачу,
\[
\|\psi_h\| \to 0,
\qquad h\to 0.
\]

Отсюда
\[
\|e_h\| \to 0,
\qquad h\to 0.
\]

Это означает сходимость.

Обратное направление в полной общей форме требует аккуратного функционально-аналитического рассмотрения. Основная идея состоит в следующем. Если схема сходится для любых допустимых данных корректной линейной задачи, то операторы перехода от данных к решению не могут иметь нормы, неограниченно растущие при $h\to0$. В противном случае существовали бы сколь угодно малые возмущения данных, которые порождали бы конечные или растущие возмущения решения, что противоречило бы сходимости для всех данных. Следовательно, семейство дискретных операторов решения равномерно ограничено, то есть схема устойчива.

Таким образом, при наличии согласованности устойчивость и сходимость эквивалентны.
\end{proof}

\begin{remark}
Теорема Лакса--Рихтмайера фундаментальна для вычислительного эксперимента: она показывает, что недостаточно построить формально точную аппроксимацию. Необходимо также обеспечить устойчивость численного процесса.
\end{remark}

\section{Пример вычислительного эксперимента: уравнение теплопроводности}

\subsection{Постановка задачи}

Рассмотрим задачу:
\[
\frac{\partial u}{\partial t}
=
a^2 \frac{\partial^2 u}{\partial x^2},
\qquad
0<x<1,\quad t>0,
\]
\[
u(0,t)=0,\qquad u(1,t)=0,
\]
\[
u(x,0)=\sin(\pi x).
\]

Точное решение известно:
\[
u(x,t)=e^{-a^2\pi^2 t}\sin(\pi x).
\]

Это удобно для проверки численного метода.

\subsection{Дискретизация}

Введём сетку:
\[
x_i=ih,\qquad h=\frac{1}{N},
\]
\[
t_n=n\tau.
\]

Явная разностная схема:
\[
u_i^{n+1}
=
u_i^n
+
r(u_{i+1}^n-2u_i^n+u_{i-1}^n),
\qquad
r=\frac{a^2\tau}{h^2}.
\]

Граничные условия:
\[
u_0^n=0,\qquad u_N^n=0.
\]

Начальное условие:
\[
u_i^0=\sin(\pi x_i).
\]

\subsection{Алгоритм явной схемы}

\begin{enumerate}
    \item Задать параметры $a$, $N$, $\tau$, $T$.
    \item Вычислить шаг $h=1/N$.
    \item Вычислить число временных слоёв $M=T/\tau$.
    \item Проверить условие устойчивости:
    \[
    r=\frac{a^2\tau}{h^2}\le \frac12.
    \]
    \item Задать начальный слой:
    \[
    u_i^0=\sin(\pi x_i).
    \]
    \item Для каждого временного слоя $n=0,1,\dots,M-1$:
    \begin{enumerate}
        \item положить $u_0^{n+1}=0$, $u_N^{n+1}=0$;
        \item для внутренних узлов $i=1,\dots,N-1$ вычислить
        \[
        u_i^{n+1}
        =
        u_i^n+r(u_{i+1}^n-2u_i^n+u_{i-1}^n).
        \]
    \end{enumerate}
    \item Сравнить численное решение с точным:
    \[
    u(x_i,T)=e^{-a^2\pi^2T}\sin(\pi x_i).
    \]
    \item Вычислить норму ошибки:
    \[
    E_h=\max_i |u_i^M-u(x_i,T)|.
    \]
\end{enumerate}

\subsection{Псевдокод}

\begin{lstlisting}[style=pseudo]
Input: a, N, tau, T
h = 1 / N
M = T / tau
r = a^2 * tau / h^2

if r > 1/2:
    print("Warning: explicit scheme may be unstable")

for i = 0..N:
    x[i] = i*h
    u[i] = sin(pi*x[i])

for n = 0..M-1:
    u_new[0] = 0
    u_new[N] = 0

    for i = 1..N-1:
        u_new[i] = u[i] + r*(u[i+1] - 2*u[i] + u[i-1])

    u = u_new

for i = 0..N:
    exact[i] = exp(-a^2*pi^2*T)*sin(pi*x[i])
    error[i] = abs(u[i] - exact[i])

E = max(error)
Output: u, E
\end{lstlisting}

\subsection{Минимальный численный пример}

Пусть:
\[
a=1,\qquad N=4,\qquad h=\frac14,
\]
\[
\tau=\frac{1}{32},\qquad r=\frac{\tau}{h^2}
=
\frac{1/32}{1/16}
=
\frac12.
\]

Узлы:
\[
x_0=0,\quad x_1=\frac14,\quad x_2=\frac12,\quad x_3=\frac34,\quad x_4=1.
\]

Начальные значения:
\[
u_0^0=0,
\]
\[
u_1^0=\sin\frac{\pi}{4}=\frac{\sqrt2}{2},
\]
\[
u_2^0=\sin\frac{\pi}{2}=1,
\]
\[
u_3^0=\sin\frac{3\pi}{4}=\frac{\sqrt2}{2},
\]
\[
u_4^0=0.
\]

Для первого временного слоя:
\[
u_i^{1}
=
u_i^0
+
\frac12(u_{i+1}^0-2u_i^0+u_{i-1}^0).
\]

Вычислим:
\[
u_1^1
=
\frac{\sqrt2}{2}
+
\frac12
\left(
1-2\frac{\sqrt2}{2}+0
\right)
=
\frac12.
\]

\[
u_2^1
=
1+
\frac12
\left(
\frac{\sqrt2}{2}-2+\frac{\sqrt2}{2}
\right)
=
1+\frac12(\sqrt2-2)
=
\frac{\sqrt2}{2}.
\]

\[
u_3^1
=
\frac{\sqrt2}{2}
+
\frac12
\left(
0-2\frac{\sqrt2}{2}+1
\right)
=
\frac12.
\]

Итак:
\[
u^1=
\left(
0,\,
\frac12,\,
\frac{\sqrt2}{2},\,
\frac12,\,
0
\right).
\]

Температура уменьшилась, что соответствует физическому смыслу задачи теплопроводности при нулевых температурах на концах стержня.

\section{Устойчивость явной схемы для уравнения теплопроводности}

\begin{theorem}
Явная разностная схема
\[
u_i^{n+1}
=
u_i^n+r(u_{i+1}^n-2u_i^n+u_{i-1}^n)
\]
для одномерного уравнения теплопроводности устойчива в норме максимума при условии
\[
0\le r\le \frac12.
\]
\end{theorem}

\begin{proof}
Запишем схему в виде:
\[
u_i^{n+1}
=
r u_{i-1}^n + (1-2r)u_i^n + r u_{i+1}^n.
\]

Если
\[
0\le r\le \frac12,
\]
то коэффициенты
\[
r,\qquad 1-2r,\qquad r
\]
неотрицательны, а их сумма равна:
\[
r+(1-2r)+r=1.
\]

Следовательно, $u_i^{n+1}$ является выпуклой комбинацией трёх значений предыдущего слоя:
\[
u_{i-1}^n,\quad u_i^n,\quad u_{i+1}^n.
\]

Отсюда:
\[
|u_i^{n+1}|
\le
r|u_{i-1}^n|+(1-2r)|u_i^n|+r|u_{i+1}^n|.
\]

Пусть
\[
\|u^n\|_\infty=\max_i |u_i^n|.
\]

Тогда:
\[
|u_i^{n+1}|
\le
r\|u^n\|_\infty+(1-2r)\|u^n\|_\infty+r\|u^n\|_\infty
=
\|u^n\|_\infty.
\]

Следовательно,
\[
\|u^{n+1}\|_\infty \le \|u^n\|_\infty.
\]

По индукции:
\[
\|u^n\|_\infty \le \|u^0\|_\infty.
\]

Это означает, что ошибки и возмущения не усиливаются от слоя к слою. Поэтому схема устойчива в норме максимума.
\end{proof}

\begin{remark}
Если $r>1/2$, коэффициент $1-2r$ становится отрицательным, и выпуклая комбинация исчезает. В этом случае возможно нарастание численных осцилляций.
\end{remark}

\section{Инструментальная поддержка вычислительного эксперимента}

\subsection{Понятие инструментальной поддержки}

\begin{definition}
\textbf{Инструментальная поддержка вычислительного эксперимента} --- совокупность аппаратных, программных, алгоритмических, информационных и организационных средств, обеспечивающих построение модели, проведение расчётов, хранение данных, анализ результатов и воспроизводимость исследования.
\end{definition}

Инструментальная поддержка включает:
\begin{itemize}
    \item вычислительную технику;
    \item операционные системы;
    \item языки программирования;
    \item библиотеки численных методов;
    \item системы компьютерной алгебры;
    \item пакеты математического моделирования;
    \item средства визуализации;
    \item базы данных;
    \item системы управления версиями;
    \item средства автоматизации расчётов;
    \item средства документирования и публикации результатов.
\end{itemize}

\subsection{Аппаратная поддержка}

К аппаратным средствам относятся:
\begin{itemize}
    \item персональные компьютеры;
    \item рабочие станции;
    \item серверы;
    \item вычислительные кластеры;
    \item суперкомпьютеры;
    \item графические процессоры;
    \item облачные вычислительные ресурсы.
\end{itemize}

Выбор аппаратной платформы зависит от:
\begin{itemize}
    \item размерности задачи;
    \item объёма данных;
    \item требуемой точности;
    \item времени расчёта;
    \item возможности параллелизации;
    \item бюджета.
\end{itemize}

\subsection{Программная поддержка}

Программные средства вычислительного эксперимента можно разделить на несколько классов.

\subsubsection{Языки программирования}

Наиболее часто используются:
\begin{itemize}
    \item Fortran --- традиционно применяется в вычислительной математике и инженерных расчётах;
    \item C и C++ --- используются для высокопроизводительных программ;
    \item Python --- широко применяется для прототипирования, анализа данных и научных вычислений;
    \item Julia --- язык, ориентированный на научные вычисления;
    \item MATLAB/Octave --- удобны для матричных вычислений и моделирования;
    \item R --- применяется в статистике и анализе данных.
\end{itemize}

\subsubsection{Библиотеки численных методов}

Типичные библиотеки:
\begin{itemize}
    \item BLAS --- базовые операции линейной алгебры;
    \item LAPACK --- решение систем линейных уравнений, задачи собственных значений;
    \item PETSc --- параллельное решение систем уравнений и задач математической физики;
    \item Trilinos --- набор библиотек для научных и инженерных вычислений;
    \item FFTW --- быстрое преобразование Фурье;
    \item NumPy и SciPy --- научные вычисления в Python.
\end{itemize}

\subsubsection{Пакеты моделирования}

К специализированным системам относятся:
\begin{itemize}
    \item ANSYS;
    \item COMSOL Multiphysics;
    \item OpenFOAM;
    \item Abaqus;
    \item Simulink;
    \item Modelica-среды;
    \item FEniCS;
    \item FreeFEM.
\end{itemize}

Они позволяют строить модели, задавать геометрию, граничные условия, сетки, выполнять расчёты и визуализировать результаты.

\subsection{Средства визуализации}

Визуализация необходима для:
\begin{itemize}
    \item обнаружения ошибок;
    \item качественного анализа решений;
    \item сравнения вариантов;
    \item представления результатов.
\end{itemize}

Используются:
\begin{itemize}
    \item графики функций;
    \item тепловые карты;
    \item контурные графики;
    \item трёхмерные поверхности;
    \item анимации;
    \item фазовые портреты;
    \item интерактивные панели.
\end{itemize}

\subsection{Средства управления данными}

Вычислительный эксперимент часто порождает большие объёмы данных:
\begin{itemize}
    \item входные параметры;
    \item сетки;
    \item промежуточные результаты;
    \item временные слои;
    \item логи работы программы;
    \item итоговые таблицы и графики.
\end{itemize}

Для их хранения применяются:
\begin{itemize}
    \item текстовые форматы;
    \item CSV;
    \item JSON;
    \item XML;
    \item HDF5;
    \item NetCDF;
    \item SQL- и NoSQL-базы данных.
\end{itemize}

\subsection{Средства воспроизводимости}

Воспроизводимость является важным требованием к вычислительному эксперименту.

Для обеспечения воспроизводимости необходимо фиксировать:
\begin{itemize}
    \item версию исходного кода;
    \item входные данные;
    \item параметры расчёта;
    \item версии библиотек;
    \item операционную систему;
    \item компилятор или интерпретатор;
    \item аппаратную платформу;
    \item случайные зерна генераторов псевдослучайных чисел;
    \item сценарии запуска.
\end{itemize}

Инструменты:
\begin{itemize}
    \item Git;
    \item Docker;
    \item Conda;
    \item Make;
    \item Snakemake;
    \item Jupyter Notebook;
    \item системы непрерывной интеграции.
\end{itemize}

\section{Типовой алгоритм вычислительного эксперимента}

\subsection{Алгоритм}

\begin{enumerate}
    \item Сформулировать цель исследования.
    \item Описать объект и выделить существенные факторы.
    \item Построить математическую модель.
    \item Определить область применимости модели.
    \item Проверить размерности и физический смысл уравнений.
    \item Выбрать численный метод.
    \item Построить дискретную модель.
    \item Оценить аппроксимацию, устойчивость и ожидаемую сходимость.
    \item Реализовать алгоритм в программе.
    \item Провести тестирование на задачах с известным решением.
    \item Выполнить серию расчётов.
    \item Провести анализ чувствительности к параметрам.
    \item Оценить погрешность.
    \item Сравнить результаты с экспериментом или теорией.
    \item Сделать выводы и при необходимости уточнить модель.
\end{enumerate}

\subsection{Псевдокод}

\begin{lstlisting}[style=pseudo]
Input: physical problem, parameters, accuracy requirements

1. Define purpose of simulation
2. Build mathematical model
3. Check assumptions and applicability
4. Choose numerical method
5. Discretize equations
6. Analyze approximation and stability
7. Implement algorithm
8. Verify code using test problems
9. Run computational experiments
10. Estimate errors
11. Validate model using experimental or reference data
12. Analyze and visualize results
13. Document code, data, parameters and conclusions

Output: numerical results, error estimates, conclusions
\end{lstlisting}

\section{Минимальный пример полного вычислительного эксперимента}

Рассмотрим задачу приближённого вычисления интеграла:
\[
I=\int_0^1 x^2\,dx.
\]

Точное значение:
\[
I=\frac13.
\]

\subsection{Математическая модель}

Интеграл можно интерпретировать как площадь под графиком функции:
\[
f(x)=x^2
\]
на отрезке $[0,1]$.

\subsection{Численный метод}

Используем формулу левых прямоугольников:
\[
I_N
=
h\sum_{i=0}^{N-1} f(x_i),
\qquad
x_i=ih,\qquad h=\frac1N.
\]

\subsection{Алгоритм}

\begin{enumerate}
    \item Задать число разбиений $N$.
    \item Вычислить $h=1/N$.
    \item Положить $S=0$.
    \item Для $i=0,\dots,N-1$:
    \[
    x_i=ih,
    \]
    \[
    S=S+x_i^2.
    \]
    \item Вычислить
    \[
    I_N=hS.
    \]
    \item Оценить ошибку:
    \[
    E_N=\left|I_N-\frac13\right|.
    \]
\end{enumerate}

\subsection{Псевдокод}

\begin{lstlisting}[style=pseudo]
Input: N
h = 1/N
S = 0

for i = 0..N-1:
    x = i*h
    S = S + x^2

I = h*S
E = abs(I - 1/3)

Output: I, E
\end{lstlisting}

\subsection{Численный расчёт при $N=4$}

\[
h=\frac14.
\]

Узлы:
\[
x_0=0,\quad x_1=\frac14,\quad x_2=\frac12,\quad x_3=\frac34.
\]

Тогда:
\[
I_4
=
\frac14
\left(
0^2+\left(\frac14\right)^2+\left(\frac12\right)^2+\left(\frac34\right)^2
\right).
\]

\[
I_4
=
\frac14
\left(
0+\frac1{16}+\frac14+\frac9{16}
\right)
=
\frac14\cdot\frac{14}{16}
=
\frac{14}{64}
=
0.21875.
\]

Точная величина:
\[
I=\frac13\approx 0.33333.
\]

Ошибка:
\[
E_4
=
\left|
0.21875-0.33333
\right|
\approx
0.11458.
\]

При увеличении $N$ ошибка уменьшается.

\section{Планирование вычислительного эксперимента}

Вычислительный эксперимент должен быть не случайным набором запусков программы, а спланированным исследованием.

\subsection{Основные элементы планирования}

\begin{itemize}
    \item выбор диапазона параметров;
    \item выбор шага изменения параметров;
    \item определение контрольных величин;
    \item выбор критериев остановки;
    \item определение метрик качества;
    \item оценка требуемых ресурсов;
    \item фиксация сценариев запуска.
\end{itemize}

\subsection{Параметрическое исследование}

Пусть результат вычислений зависит от параметра $\alpha$:
\[
y = F(\alpha).
\]

Тогда проводится серия расчётов:
\[
\alpha_1,\alpha_2,\dots,\alpha_m,
\]
и строится зависимость:
\[
y_j=F(\alpha_j).
\]

Целью может быть:
\begin{itemize}
    \item поиск критического значения параметра;
    \item выявление устойчивых режимов;
    \item обнаружение бифуркаций;
    \item оптимизация результата.
\end{itemize}

\subsection{Анализ чувствительности}

\begin{definition}
\textbf{Анализ чувствительности} --- исследование изменения результата при изменении входных параметров.
\end{definition}

Если
\[
y=F(p),
\]
то локальная чувствительность может быть оценена производной:
\[
S(p)=\frac{dF}{dp}.
\]

В относительном виде:
\[
S_{\text{rel}}(p)
=
\frac{p}{F(p)}\frac{dF}{dp}.
\]

Численно производную можно приблизить конечной разностью:
\[
\frac{dF}{dp}
\approx
\frac{F(p+\Delta p)-F(p)}{\Delta p}.
\]

\section{Верификация программной реализации}

\subsection{Основные способы верификации}

\begin{enumerate}
    \item \textbf{Сравнение с точным решением.}

    Если известно аналитическое решение, численный результат сравнивается с ним.

    \item \textbf{Метод сгущения сетки.}

    Расчёты проводятся на нескольких сетках:
    \[
    h,\quad \frac h2,\quad \frac h4.
    \]
    Если метод имеет порядок $p$, ошибка должна убывать примерно как $h^p$.

    \item \textbf{Проверка законов сохранения.}

    Например, для механических систем проверяется сохранение энергии, импульса или массы.

    \item \textbf{Сравнение с эталонным кодом.}

    Результаты сравниваются с известной проверенной программой.

    \item \textbf{Тестирование предельных случаев.}

    Проверяются ситуации, для которых результат заранее понятен.
\end{enumerate}

\subsection{Оценка порядка сходимости}

Пусть ошибки на сетках $h$ и $h/2$ равны:
\[
E_h \approx C h^p,
\]
\[
E_{h/2} \approx C\left(\frac h2\right)^p.
\]

Тогда:
\[
\frac{E_h}{E_{h/2}}
\approx
2^p.
\]

Следовательно:
\[
p \approx \log_2\frac{E_h}{E_{h/2}}.
\]

\begin{example}
Если при уменьшении шага в два раза ошибка уменьшается примерно в четыре раза:
\[
\frac{E_h}{E_{h/2}}\approx 4,
\]
то
\[
p\approx \log_2 4 = 2.
\]
Следовательно, метод имеет второй порядок сходимости.
\end{example}

\section{Валидация модели}

Валидация отвечает на вопрос, насколько хорошо модель описывает реальность.

Методы валидации:
\begin{itemize}
    \item сравнение с экспериментальными данными;
    \item сравнение с наблюдениями;
    \item экспертная оценка;
    \item проверка предсказательной способности;
    \item статистическая оценка отклонений;
    \item кросс-валидация для моделей данных.
\end{itemize}

\subsection{Критерии согласия}

Пусть имеются экспериментальные значения $y_i$ и расчётные значения $\widehat{y}_i$. Часто используются следующие метрики.

Средняя абсолютная ошибка:
\[
MAE
=
\frac1n\sum_{i=1}^n |y_i-\widehat{y}_i|.
\]

Среднеквадратичная ошибка:
\[
MSE
=
\frac1n\sum_{i=1}^n (y_i-\widehat{y}_i)^2.
\]

Корень из среднеквадратичной ошибки:
\[
RMSE
=
\sqrt{
\frac1n\sum_{i=1}^n (y_i-\widehat{y}_i)^2
}.
\]

Относительная ошибка:
\[
\varepsilon_i
=
\frac{|y_i-\widehat{y}_i|}{|y_i|}.
\]

\section{Роль суперкомпьютеров и параллельных вычислений}

Многие вычислительные эксперименты требуют больших вычислительных ресурсов. Например:
\begin{itemize}
    \item моделирование климата;
    \item расчёт турбулентных течений;
    \item квантово-химические вычисления;
    \item моделирование распространения волн;
    \item анализ больших данных;
    \item обучение крупных моделей машинного обучения.
\end{itemize}

\subsection{Параллелизация}

\begin{definition}
\textbf{Параллелизация} --- разбиение вычислительной задачи на части, которые могут выполняться одновременно.
\end{definition}

Основные виды параллелизма:
\begin{itemize}
    \item параллелизм по данным;
    \item параллелизм по задачам;
    \item конвейерный параллелизм;
    \item гибридный параллелизм.
\end{itemize}

\subsection{Закон Амдала}

\begin{theorem}[Закон Амдала]
Пусть доля вычислений, допускающая параллельное выполнение, равна $p$, а оставшаяся доля $1-p$ выполняется последовательно. Тогда максимальное ускорение при использовании $N$ процессоров равно
\[
S_N
=
\frac{1}{(1-p)+\frac{p}{N}}.
\]
\end{theorem}

\begin{proof}
Пусть время выполнения программы на одном процессоре равно $T_1$.

Последовательная часть занимает:
\[
T_{\text{seq}}=(1-p)T_1.
\]

Параллельная часть занимает:
\[
T_{\text{par}}=pT_1.
\]

Если параллельная часть идеально распределяется между $N$ процессорами, то её время выполнения становится:
\[
\frac{pT_1}{N}.
\]

Общее время на $N$ процессорах:
\[
T_N=(1-p)T_1+\frac{pT_1}{N}.
\]

Ускорение определяется как:
\[
S_N=\frac{T_1}{T_N}.
\]

Подставляя выражение для $T_N$, получаем:
\[
S_N
=
\frac{T_1}{(1-p)T_1+\frac{pT_1}{N}}
=
\frac{1}{(1-p)+\frac{p}{N}}.
\]

Теорема доказана.
\end{proof}

\begin{remark}
Закон Амдала показывает, что даже при бесконечном числе процессоров ускорение ограничено:
\[
\lim_{N\to\infty} S_N = \frac{1}{1-p}.
\]
Если $p=0.9$, то максимальное ускорение не превосходит $10$.
\end{remark}

\section{Качество вычислительного эксперимента}

Качественный вычислительный эксперимент должен удовлетворять следующим требованиям:
\begin{itemize}
    \item математическая модель должна быть адекватной цели исследования;
    \item задача должна быть корректно поставленной или регуляризованной;
    \item численный метод должен быть устойчивым и сходящимся;
    \item программа должна быть проверена;
    \item погрешности должны быть оценены;
    \item результаты должны быть воспроизводимы;
    \item выводы должны учитывать ограничения модели.
\end{itemize}

\subsection{Типичные ошибки при проведении вычислительных экспериментов}

\begin{enumerate}
    \item Использование модели вне области её применимости.
    \item Отсутствие проверки размерностей.
    \item Игнорирование устойчивости численной схемы.
    \item Недостаточное сгущение сетки.
    \item Сравнение результатов без оценки погрешности.
    \item Использование случайных данных без фиксации зерна генератора.
    \item Отсутствие контроля версий.
    \item Неправильная визуализация.
    \item Подмена валидации красивой картинкой.
    \item Неполное документирование параметров расчёта.
\end{enumerate}

\section{Заключение}

Вычислительный эксперимент представляет собой систематическое исследование математической модели с помощью численных методов и компьютерных средств. Его основой служит математическое моделирование, а центральными требованиями являются корректность постановки, аппроксимация, устойчивость, сходимость, верификация, валидация и воспроизводимость.

Инструментальная поддержка вычислительного эксперимента включает аппаратные ресурсы, языки программирования, библиотеки численных методов, системы моделирования, средства визуализации, управления данными и обеспечения воспроизводимости. Современный вычислительный эксперимент является не просто запуском программы, а полным научным циклом: от постановки задачи и построения модели до анализа ошибок, интерпретации результатов и документирования исследования.

\end{document}
